\documentclass[12pt]{article}

\usepackage[a4paper, total={6in, 8in}]{geometry}
\usepackage{textcomp}

\begin{document}
\setlength{\parindent}{0pt}

\def \t {\textrightarrow}
\def \v {\vspace{1em}}
\def \bi {\begin{itemize}}
\def \ei {\end{itemize}}
\def \s[#1] {\section*{#1}}
\def \ss[#1] {\subsection*{#1}}
\def \sss[#1] {\subsubsection*{#1}}

\s[Uno nessuno centomila]
La parabola inizia già con il Fu mattia pascal \\
Un'esperienza di vita gli dimostra che non si puo afferrare una certezza anche minina a Priandello \\
Società irrazionalistica, nulla ha una spiegazione \t l'individuo soffocato dalla società lo è in primis dal sistema familiare, in tutti e tre i romanzi \\
Il Moscarda trova nella usa domanda di senso la molteplicità delle ripsorte \t quando chiede alla moglie, riceve una ripsota che non si darebbe \\
In Pascal è invece una selva \\
Uno contro la società, contro la visione critica della società che non va oltre la cristallizazione della forma, e perde di vista il flusso continuo \\
Finiamo per essere personaggi \t realtà sociale complessa, nella quale si mette in gioco il concetto di identità e unità della coscienza \\
La psicologina non nasce con Freud \t Alfred Vine pubblica "Le alterazioni della personalità" \t sostiene che in ciascuno conviverebbero molteplici personalità, delle quali siamo piu o meno consapevoli \\
Noi non siamo univoci nel dimostrare noi stessi \t questo anticipa il mondo di Freud, che Pirandello conosce superficialmente e non apprezza in toto \\
L'identità sociale non corrisponde a quella di noi stessi \t noi siamo inafferrabili anche a noi stessi \\
Anticipazione delle intelligenze multiple, e della dissociazione \\
L'idea stessa della integrità è in gioco \t Pirandello vuole salvare almeno la coscienza \t mentre svevo non la salverà neanche piu \\
La parola oggettività è bandita, il capirsi è impossibile, il paradigma è l'incomunicabilità, l'alienazione = negazione di ongi processo di conoscenza oggettiva, il punto di vista gnoseologico è calzante nella definizione di Grosserman \\
Qua si coglie la quintessenza della filosofia pirandelliana, che è a se stante e non sottosta a nessuna corrente filosofica

\v

Nella fagocitata successione di eventi che caratterizzano l'opera del moscarda, è come un insieme di burattini che non trovano il loro ruolo \\
Questo si incarna nella meccanicità d'azione di alcuni personaggi \t ma se il titolo dell'opera è espressivo, l apice dell evoluzione di una coscienza si ha con il fu mattia pascal \\
Da li ci sono i fallimentari tentativi di trovare un posto sulla terra, e di non essere al centro di nulla se non la tragedia del vivere \\
Poi la forma in cui siamo incastrati \t una deviazione da quella scardina il complesso della nostra esistenza \t non esistiamo + se non veniamo ritrovati nella forma \\
La forma ci rende inesistenti \t questo è guardato attraverso l'avvertimento del contrario \t e la riflessione del sentimento del contraio che dice che è il caso, su cui non abbiamo potere \\
È il caso che ci fa vederci vivere, siamo forestieri della vita

\v

Questo si legge nella complessità della colloquialità delle opere pirandelliane \t un crescendo che sfocia fino ad essere teatrale quasi, come dei marchinegni pilotati a scatti \\
Il linguaggio in questo segue la quotidianità e dimostra la incomunicabilità \t il parlare e incepparsi, il parale a ruota libera, è lo specchio del conflitto interiore che vivono dentro di se \\
Questo io che è destituito completamente \t non possiamo vivere senza la consapevolezza dell'io, e la sua decostituzione porta a chiedersi chi siamo \\
La riflessione di questo vedersi vivere si incardina nelle opere attraverso quello che si chiama il ragioneur, quello che continua a congetturare \\
Ci mostra un continuo discutere dei peronsaggi con se stessi \t l'arte che non deriva dal classicismo, è filologo, la difficoltà del vivere fa capire lo scollamento dalle certezze che il mondo romantico aveva instillato nell'uomo \\
Ne esce una modernità \t le pulsioni vitali si scontrano con la forma, la vita è impossibile dietro a una maschera \\
Tutto cio che è sociale è un apparato troppo lontano per l uomo, quasi kafkiano, in cui viene distrutta la persona e non il personaggio \\
Cosi noi non possiamo vivere e siamo costretti a guardarci vivere \t ci autoninganniamo quando assorbiamo la maschera che gli altri ci hanno dato \t se vogliamo essere persona rischiamo di non esistere \\
E questo è cio che caratterizza tutte le opere

\v

Marionette tirate da un filo \t il volto della vita tratteggiata da pirandello \\
I personaggi sono molto attivi nel dibattere tra loro \t questo li rende a un passo dal mondo teatrale \\
Teatrato dionesco, la cantratice calva, la sedia \t il grottesco \t Bertold Brecht, l'assurdo \\
Mistero enigmatico della vita, il teatro nel teatro \t dal 21 rappresentazione di "6 personaggi in cerca d'autore": il mondo pirandelliano è ben presente \\
Il monoatto dall "Uomo dal fiore in bocca" e l'"Enrico IV", il senso della vita quando si sta perdendo, la ricerca del proprio posto nel mondo e l'impossibilita di trovarlo (6 personaggi) \\
Enrivo IV \t impossibilita del sopravvivere \\
Il teatro prosegue anche dal 28 al 30 (oltre quella del 21-22) \t e parti metateatrali sono presenti anche nel romanzo incompiuto "I giganti della montagna" \\
Disquisizione teatro nel teatro, gioco della forma e della vita, della autonomia dell'autore = regista \\
Visione del teatro metatratale, e il denominatore è che come il dissidio dell'io è presente nell autore e nei personaggi, qui il dissidio è maschera e persona (neanche personaggio)
\end{document}
